\begin{exsol}{6}{Varian Ch.~32, Review Question~3}{628}{Foundations}{B1}
\asked Decide whether knowing the contract curve of an exchange economy is enough to know where
trade between the two consumers will end up, and justify the verdict.

\method Pareto efficiency in the Edgeworth box and the contract curve (Varian \S32.3,
pp.~605--607); the role of the endowment and of the trading mechanism (competitive equilibrium,
take-it-or-leave-it offers, Nash bargaining); the course notes, Part~II, sections ``The
question: is bargaining indeterminate?'' and ``Dividing the gains: the Nash bargaining solution''.

\solution
\textbf{Verdict: false.} The contract curve tells us \emph{which kind} of point trade ends at, not
\emph{which} point.

\emph{(a) What the contract curve is.} In a two-person, two-good exchange economy with total
endowment $(\bar\omega_1,\bar\omega_2)$, an allocation is Pareto efficient when no reallocation
makes one consumer better off without making the other worse off. At an interior efficient
allocation the indifference curves are tangent,
\[
  MRS^A_{12}(x^A) \;=\; MRS^B_{12}(\bar\omega - x^A),
\]
because crossing curves leave a lens of mutually preferred allocations (Varian p.~606). The
contract curve (Pareto set) is the locus of all such points; it typically runs from $A$'s origin
to $B$'s origin, and it is defined by preferences and totals alone.

\emph{(b) What it does tell us.} If trade is voluntary and the parties keep trading until no
mutually advantageous trade remains, the final allocation is on the contract curve: a point off
the curve would not be a ``final contract'' (Varian p.~606). Voluntary trade also requires that
each consumer end up at least as well off as at the endowment $\omega$. So trade ends in the
segment of the contract curve inside the lens through $\omega$ (Varian p.~607; this segment is the
core of the two-person economy). This is exactly Varian's answer (p.~A27): trade ends ``somewhere on
the curve; however, we don't know where.'' We agree, and add what decides the point.

\emph{(c) What it does not tell us.} Two things, neither of which is in the contract curve:
\begin{enumerate}[label=(\roman*)]
  \item \emph{The endowment.} A different $\omega$ selects a different segment.
  \item \emph{The mechanism.} Within the segment, the point is set by how trade happens: a
  competitive (Walrasian) equilibrium, a take-it-or-leave-it offer by one side, or a bargaining
  rule that splits the gains according to bargaining power and patience.
\end{enumerate}

\emph{Example.} Let $u_A = \ln x_1^A + \ln x_2^A$, $u_B = \ln x_1^B + \ln x_2^B$, totals
$(10,10)$ and endowment $\omega^A = (8,2)$, $\omega^B = (2,8)$. Tangency gives
$x_2^A/x_1^A = (10-x_2^A)/(10-x_1^A)$, i.e.\ the contract curve is the diagonal
$x_1^A = x_2^A = t$. Individual rationality on the diagonal:
\[
  u_A:\; t^2 \ge 8\cdot 2 = 16 \;\Rightarrow\; t \ge 4, \qquad
  u_B:\; (10-t)^2 \ge 16 \;\Rightarrow\; t \le 6 .
\]
So trade ends somewhere on the segment $t \in [4,6]$, and the mechanism picks the point:
\begin{center}
\begin{tabular}{lc}
\toprule
Mechanism & Final allocation of $A$ \\
\midrule
$A$ makes a take-it-or-leave-it offer & $t = 6$ ($B$ indifferent to $\omega$) \\
$B$ makes a take-it-or-leave-it offer & $t = 4$ ($A$ indifferent to $\omega$) \\
Competitive equilibrium ($p_1/p_2 = 1$, $A$'s income $10$) & $t = 5$ \\
Symmetric Nash bargaining & $t = 5$ \\
\bottomrule
\end{tabular}
\end{center}
All four outcomes are on the same contract curve; they differ in who gets the gains from trade.

\begin{center}
\begin{tikzpicture}[scale=0.5,>=Stealth]
  \draw (0,0) rectangle (10,10);
  \node[below left] at (0,0) {$O_A$};
  \node[above right] at (10,10) {$O_B$};
  \node[below] at (5,-0.2) {\footnotesize good 1 ($A$)};
  \node[left, rotate=90, anchor=south] at (-0.3,5) {\footnotesize good 2 ($A$)};
  \draw[thick, hdr] (0,0) -- (10,10);
  \node[hdr, font=\footnotesize, below right] at (7.4,7.4) {contract curve};
  \draw[gray] plot[domain=1.6:10, samples=60] (\x,{16/\x});
  \draw[gray, dashed] plot[domain=0:8.4, samples=60] (\x,{10-16/(10-\x)});
  \draw[ultra thick, ans] (4,4) -- (6,6);
  \fill (8,2) circle (3pt) node[right] {$\omega$};
  \fill (5,5) circle (3pt) node[right] {\footnotesize CE / Nash};
  \fill (4,4) circle (2pt); \fill (6,6) circle (2pt);
\end{tikzpicture}
\end{center}
{\footnotesize Solid grey: $A$'s indifference curve through $\omega$; dashed: $B$'s. Green: the
part of the contract curve both prefer to $\omega$, where voluntary efficient trade ends.}

\emph{(d) A caveat in the other direction.} Even ``somewhere on the curve'' presumes that trade
exhausts the gains. With private information about valuations, no bargaining procedure may reach
the contract curve (\MR{} Ch.~5, pp.~141--143; notes, Part~II, ``Private information''). The contract curve is
then the set of \emph{possible} efficient outcomes, not a prediction.

\interp This is the exchange-economy version of \MR's separation (pp.~36--38): efficiency
restricts \emph{what} is done, while bargaining power, patience and the initial position decide
\emph{who gets what}. Efficiency is a property of the outcome; the division needs a theory of the
bargaining process (Nash, Rubinstein), which is why the course moves on to bargaining after the
value maximization principle.

\verify At $t=4$, $u_A = 2\ln 4 = \ln 16 = u_A(\omega)$; at $t=6$, $u_B = 2\ln 4 = u_B(\omega)$.
With $p_1 = p_2 = 1$, $A$'s income is $8+2 = 10$ and Cobb--Douglas demand spends half on each good:
$x^A = (5,5)$, $x^B = (5,5)$, markets clear. Symmetry of the Nash product
$(2\ln t - \ln 16)(2\ln(10-t) - \ln 16)$ around $t=5$ gives $t=5$ (\texttt{check\_C.py}).

\answer{False. The contract curve is the set of efficient allocations; voluntary efficient trade
ends on its segment inside the lens through the endowment, but the point within that segment
depends on the endowment and on the trading mechanism (in the example, anywhere in $t\in[4,6]$:
$4$ or $6$ under take-it-or-leave-it offers, $5$ under competition or symmetric Nash).}
\end{exsol}
